\chapter{1977年福建卷（理）}

\section{解答题}

\begin{problem}
\begin{enumerate}
    \item 计算：\(5-3\times\left[\left(-3\frac{3}{8}\right)^{-\frac{1}{3}}+1031\times\left(0.25-2^{-2}\right)\right]\div9^0\).

    \item \(y=\frac{\cos160^\circ-\cos170^\circ}{\tan155^\circ}\) 的值是正的还是负的？为什么？

    \item 求函数 \(y=\frac{\lg\left(2-x\right)}{\sqrt{x-1}}\) 的定义域.

    \item 如图，在梯形 \(ABCD\) 中，\(DM=MP=PA\)，\(MN\parallel PQ\parallel AB\)，\(DC=2\mathrm{~cm}\)，\(AB=3.5\mathrm{~cm}\)，求 \(MN\) 和 \(PQ\) 的长.

    \bitmapfigure[max width=5.8cm]{img/1977/fujian_science/q1_fig1.png}

    \item 已知 \(\lg3=0.4771\)，\(\lg x=-3.5229\)，求 \(x\).

    \item 求 \(\displaystyle\lim_{x\to1}\frac{x-1}{x^2-3x+2}\).

    \item 解方程：\(\sqrt{4x+1}-2x+1=0\).

    \item 化简：\(\frac{a^{2n+1}-6a^{2n}+9a^{2n-1}}{a^{n+1}-4a^n+3a^{n-1}}\).

    \item 求函数 \(y=2-5x-3x^2\) 的极值.

    \item 画出下面 \(V\) 形铁块的三视图（只要画草图）.
\end{enumerate}

\bitmapfigure[max width=5.8cm]{img/1977/fujian_science/q1_fig2.png}
\end{problem}

\begin{solution}
\begin{enumerate}
    \item 因为 \(-3\frac{3}{8}=-\frac{27}{8}\)，所以 \(\left(-3\frac{3}{8}\right)^{-\frac{1}{3}}=\left(-\frac{3}{2}\right)^{-1}=-\frac{2}{3}\)，且 \(0.25-2^{-2}=\frac{1}{4}-\frac{1}{4}=0\)，\(9^0=1\). 因此原式 \(=5-3\times\left(-\frac{2}{3}\right)=7\).

    \item 在 \([0^\circ,180^\circ]\) 上，余弦函数单调递减，且 \(160^\circ<170^\circ\)，所以 \(\cos160^\circ-\cos170^\circ>0\). 又 \(155^\circ\) 在第二象限，\(\tan155^\circ<0\). 因此分子为正、分母为负，故 \(y<0\)，这个值是负的.

    \item 由对数有意义，得 \(2-x>0\)，即 \(x<2\). 由于根式在分母中，必须有 \(x-1>0\)，即 \(x>1\). 两个条件同时成立，所以定义域为 \(\left(1,2\right)\).

    \item 由相似三角形或平行线分线段成比例定理，平行截线的长度随左腰上的分点位置作线性变化. \(M\) 把 \(DA\) 分成三等份，故
    \[
    MN=DC+\frac{1}{3}\left(AB-DC\right)=2+\frac{1}{3}\left(3.5-2\right)=2.5\mathrm{~cm}
    \].
    \(P\) 到 \(D\) 的距离占 \(DA\) 的 \(\frac{2}{3}\)，故
    \[
    PQ=DC+\frac{2}{3}\left(AB-DC\right)=2+\frac{2}{3}\left(3.5-2\right)=3\mathrm{~cm}
    \].

    \item 因为 \(\lg10^4=4\)，所以
    \[
    \lg x=-3.5229=0.4771-4=\lg3-\lg10^4=\lg\frac{3}{10^4}
    \].
    对数函数是单调函数，故 \(x=\frac{3}{10^4}=0.0003\).

    \item 当 \(x\ne1,2\) 时，\(x^2-3x+2=(x-1)(x-2)\)，从而
    \[
    \frac{x-1}{x^2-3x+2}=\frac{1}{x-2}
    \].
    因此
    \[
    \displaystyle\lim_{x\to1}\frac{x-1}{x^2-3x+2}=\lim_{x\to1}\frac{1}{x-2}=-1
    \].

    \item 由原方程得 \(\sqrt{4x+1}=2x-1\)，所以必须有 \(2x-1\geqslant0\)，即 \(x\geqslant\frac{1}{2}\). 两边平方，得
    \[
    4x+1=(2x-1)^2=4x^2-4x+1
    \],
    即 \(4x(x-2)=0\)，所以 \(x=0\) 或 \(x=2\). 由 \(x\geqslant\frac{1}{2}\) 舍去 \(x=0\)，代入原方程检验 \(x=2\) 成立，故原方程的解为 \(x=2\).

    \item 分子、分母分别因式分解，得
    \[
    a^{2n+1}-6a^{2n}+9a^{2n-1}=a^{2n-1}(a-3)^2
    \],
    \[
    a^{n+1}-4a^n+3a^{n-1}=a^{n-1}(a-1)(a-3)
    \].
    在原式有意义的条件下约去公因式，得
    \[
    \frac{a^{2n+1}-6a^{2n}+9a^{2n-1}}{a^{n+1}-4a^n+3a^{n-1}}=\frac{a^n(a-3)}{a-1}
    \].
    若 \(n\geqslant2\)，原式还要求 \(a\ne0\)，并且始终要求 \(a\ne1,3\)；若 \(n=1\)，原式的定义域只要求 \(a\ne1,3\).

    \item 配方得
    \[
    y=2-5x-3x^2=\frac{49}{12}-3\left(x+\frac{5}{6}\right)^2
    \].
    因为二次项系数为负，所以函数在 \(x=-\frac{5}{6}\) 时取得最大值 \(\frac{49}{12}\)，没有最小值.
\end{enumerate}
\end{solution}

\begin{problem}
\begin{enumerate}
    \item 解不等式：\(\frac{x^2-x-6}{x^2+2x+2}<0\).

    \item 证明：\(\frac{2\cos\theta-\sin2\theta}{2\cos\theta+\sin2\theta}=\tan^2\left(\frac{90^\circ-\theta}{2}\right)\).

    \item 某中学革命师生自己动手油漆一个直径为 \(1.2\mathrm{~m}\) 的地球仪，如果每平方米面积需要油漆 \(150\mathrm{~g}\)，问共需油漆多少克？（答案保留整数）

    \item 某农机厂开展“工业学大庆”运动，在十月份生产拖拉机 \(1000\) 台. 这样，一月至十月的产量恰好完成全年生产任务. 工人同志为了加速农业机械化，计划在年底前再生产 \(2310\) 台，求十一月、十二月份平均每月增长率.
\end{enumerate}
\end{problem}

\begin{solution}
\begin{enumerate}
    \item 因为 \(x^2+2x+2=(x+1)^2+1>0\)，所以原不等式等价于
    \[
    x^2-x-6=(x-3)(x+2)<0
    \].
    因此不等式的解集为 \(\left(-2,3\right)\).

    \item 原式有意义时 \(\cos\theta\ne0\)，先化简左端：
    \[
    \frac{2\cos\theta-\sin2\theta}{2\cos\theta+\sin2\theta}=\frac{2\cos\theta-2\sin\theta\cos\theta}{2\cos\theta+2\sin\theta\cos\theta}=\frac{1-\sin\theta}{1+\sin\theta}
    \].
    另一方面，令 \(u=90^\circ-\theta\)，由半角公式得
    \[
    \tan^2\left(\frac{90^\circ-\theta}{2}\right)=\frac{1-\cos u}{1+\cos u}=\frac{1-\sin\theta}{1+\sin\theta}
    \].
    两端相等，恒等式得证. 上述变形在原式有意义时成立.

    \item 地球仪可看作半径为 \(0.6\mathrm{~m}\) 的球，其表面积为
    \[
    4\pi\times\left(0.6\right)^2=1.44\pi\mathrm{~m}^2
    \].
    所需油漆质量为 \(150\times1.44\pi=216\pi\mathrm{~g}\approx678.6\mathrm{~g}\)，保留整数得约 \(679\mathrm{~g}\).

    \item 设十一月、十二月份平均每月增长率为 \(r\)，则十一月和十二月的产量分别为 \(1000(1+r)\) 台和 \(1000(1+r)^2\) 台. 由题意
    \[
    1000(1+r)+1000(1+r)^2=2310
    \].
    令 \(q=1+r\)，则 \(q>0\)，且
    \[
    q^2+q-2.31=0
    \],
    所以 \((q-1.1)(q+2.1)=0\). 由 \(q>0\) 得 \(q=1.1\)，故 \(r=0.1=10\%\).
\end{enumerate}
\end{solution}

\begin{problem}
在半径为 \(R\) 的圆内接正六边形内，依次连结各边的中点，得一正六边形，又在这一正六边形内，再依次连结各边的中点，又得一正六边形，这样无限地继续下去，求：

\begin{enumerate}
    \item 前 \(n\) 个正六边形的周长之和 \(S_n\)；
    \item 所有这些正六边形的周长之和 \(S\).
\end{enumerate}
\end{problem}

\begin{solution}
设第一个正六边形的边长为 \(l_1\). 正六边形内接于半径为 \(R\) 的圆，所以 \(l_1=R\)，其周长为 \(6R\). 对于边长为 \(l_k\) 的正六边形，连接各边中点所得的正六边形的外接圆半径为 \(l_k\cos30^\circ=\frac{\sqrt{3}}{2}l_k\)，而正六边形的边长等于外接圆半径，故
\[
l_{k+1}=\frac{\sqrt{3}}{2}l_k
\].
因此各正六边形的周长构成首项为 \(6R\)、公比为 \(\frac{\sqrt{3}}{2}\) 的等比数列. 从而
\[
S_n=6R\frac{1-\left(\frac{\sqrt{3}}{2}\right)^n}{1-\frac{\sqrt{3}}{2}}=12R\left(2+\sqrt{3}\right)\left[1-\left(\frac{\sqrt{3}}{2}\right)^n\right]
\].
因为 \(0<\frac{\sqrt{3}}{2}<1\)，所以
\[
S=\lim_{n\to\infty}S_n=\frac{6R}{1-\frac{\sqrt{3}}{2}}=12R\left(2+\sqrt{3}\right)
\].
\end{solution}

\begin{problem}
动点 \(P(x,y)\) 到两定点 \(A(-3,0)\) 和 \(B(3,0)\) 的距离的比等于 \(2\)，求动点 \(P\) 的轨迹方程，并说明这轨迹是什么图形.
\end{problem}

\begin{solution}
按题目中定点的书写顺序取 \(\frac{PA}{PB}=2\)，则
\[
\sqrt{(x+3)^2+y^2}=2\sqrt{(x-3)^2+y^2}
\].
两边平方并整理，得
\[
x^2+6x+9+y^2=4x^2-24x+36+4y^2
\],
即 \(x^2-10x+9+y^2=0\). 配方得
\[
\left(x-5\right)^2+y^2=16
\].
因此轨迹是圆心为 \((5,0)\)、半径为 \(4\) 的圆.
\end{solution}

\begin{problem}
某大队在农田基本建设的规划中，要测定被障碍物隔开的两点 \(A\)，\(P\) 之间的距离，他们土法上马，在障碍物的两侧，选取两点 \(B\) 和 \(C\)（如图），测得 \(AB=AC=50\mathrm{~m}\)，\(\angle BAC=60^\circ\)，\(\angle ABP=120^\circ\)，\(\angle ACP=135^\circ\)，求 \(A\) 和 \(P\) 之间的距离.（答案可用最简根式表示）

\bitmapfigure[max width=6.4cm]{img/1977/fujian_science/q5_fig1.png}
\end{problem}

\begin{solution}
由 \(AB=AC\) 且 \(\angle BAC=60^\circ\)，知 \(\triangle ABC\) 为等边三角形. 设 \(\beta=\angle BAP\). 按图中位置，\(\angle PAC=60^\circ-\beta\). 在 \(\triangle ABP\) 中，
\[
\angle APB=180^\circ-120^\circ-\beta=60^\circ-\beta
\].
在 \(\triangle ACP\) 中，
\[
\angle APC=180^\circ-135^\circ-\left(60^\circ-\beta\right)=\beta-15^\circ
\].
由正弦定理，分别在两个三角形中得到
\[
AP=\frac{50\sin120^\circ}{\sin\left(60^\circ-\beta\right)}=\frac{25\sqrt{3}}{\sin\left(60^\circ-\beta\right)}
\],
\[
AP=\frac{50\sin135^\circ}{\sin\left(\beta-15^\circ\right)}=\frac{25\sqrt{2}}{\sin\left(\beta-15^\circ\right)}
\].
    由三角形内角均为正，得 \(15^\circ<\beta<60^\circ\). 令 \(u=\beta-15^\circ\)，则 \(0^\circ<u<45^\circ\)，且 \(60^\circ-\beta=45^\circ-u\)，从而
\[
\sqrt{3}\sin u=\sqrt{2}\sin\left(45^\circ-u\right)=\cos u-\sin u
\].
因此 \((1+\sqrt{3})\sin u=\cos u\)，即
\[
\tan u=\frac{1}{1+\sqrt{3}}=\frac{\sqrt{3}-1}{2}
\].
于是
\[
AP^2=\frac{(25\sqrt{2})^2}{\sin^2u}=1250\frac{1+\tan^2u}{\tan^2u}
\].
因为 \(\tan^2u=1-\frac{\sqrt{3}}{2}\)，所以
\(\frac{1+\tan^2u}{\tan^2u}=\frac{4-\sqrt{3}}{2-\sqrt{3}}=5+2\sqrt{3}\).
所以
\[
AP^2=1250\left(5+2\sqrt{3}\right)=625\left(10+4\sqrt{3}\right)
\].
由距离为正，得 \(AP=25\sqrt{10+4\sqrt{3}}\mathrm{~m}\).
\end{solution}

\begin{problem}
已知双曲线 \(\frac{x^2}{24\alpha}-\frac{y^2}{16\cot\alpha}=1\)（\(\alpha\) 为锐角）和圆 \(\left(x-m\right)^2+y^2=r^2\) 相切于点 \(A\left(4\sqrt{3},4\right)\)，求 \(\alpha\)，\(m\)，\(r\) 的值.
\end{problem}

\begin{solution}
按题面未标注角度单位，以下按弧度制理解. 将点 \(A\) 代入双曲线方程，得
\[
\frac{48}{24\alpha}-\frac{16}{16\cot\alpha}=1
\],
即
\[
\frac{2}{\alpha}-\tan\alpha=1
\].
    令 \(f(\alpha)=\frac{2}{\alpha}-\tan\alpha-1\). 在 \(\left(0,\frac{\pi}{2}\right)\) 内，\(f'(\alpha)=-\frac{2}{\alpha^2}-\sec^2\alpha<0\)，且 \(f(\alpha)\) 从正无穷减至负无穷，所以该方程有唯一的锐角解. 例如 \(f(0.8924934)>0\)，\(f(0.8924935)<0\)，用二分法逐次逼近得
\[
\alpha\approx0.8924934185
\].

对双曲线方程作隐式求导，得
\[
\frac{x}{12\alpha}-\frac{yy'}{8\cot\alpha}=0
\].
所以 \(y'=\frac{2x\cot\alpha}{3\alpha y}\).
故双曲线在 \(A\) 点的切线斜率为
\[
k=\frac{2\sqrt{3}\cot\alpha}{3\alpha}
\].
圆的方程求导得 \(y'= -\frac{x-m}{y}=\frac{m-x}{y}\)，所以圆在 \(A\) 点的切线斜率为 \(\frac{m-4\sqrt{3}}{4}\). 相切意味着两斜率相等，因而
\[
m=4\sqrt{3}+\frac{8\sqrt{3}\cot\alpha}{3\alpha}\approx11.09865443
\].
最后由点 \(A\) 在圆上，得
\[
r^2=\left(4\sqrt{3}-m\right)^2+4^2
\],
从而
\[
r=\sqrt{\left(4\sqrt{3}-m\right)^2+16}\approx5.77863853
\].
因此，按题面未标注角度单位时的通常约定取弧度制，\(\alpha\approx0.8924934185\)，\(m\approx11.09865443\)，\(r\approx5.77863853\).
\end{solution}

\begin{problem}
设数列 \(1\)，\(2\)，\(4\)，\(\cdots\) 前 \(n\) 项和是 \(S_n=a+bn+cn^2+dn^3\)，求这数列的通项 \(a_n\) 的公式，并确定 \(a\)，\(b\)，\(c\)，\(d\) 的值.
\end{problem}

\begin{solution}
由前几项可知 \(S_0=0\)，\(S_1=1\)，\(S_2=1+2=3\)，\(S_3=1+2+4=7\). 代入 \(S_n=a+bn+cn^2+dn^3\)，得
\[
\begin{cases}
a=0,\\
a+b+c+d=1,\\
a+2b+4c+8d=3,\\
a+3b+9c+27d=7.
\end{cases}
\]
由此解得 \(a=0\)，\(b+c+d=1\)，\(2c+6d=1\)，\(6c+24d=4\). 所以
\(a=0\)，\(b=\frac{5}{6}\)，\(c=0\)，\(d=\frac{1}{6}\).
于是 \(S_n=\frac{5n+n^3}{6}\). 对 \(n\geqslant1\)，有
\[
a_n=S_n-S_{n-1}=\frac{n^3+5n-(n-1)^3-5(n-1)}{6}=\frac{n^2-n+2}{2}
\].
故数列的通项公式为 \(a_n=\frac{n^2-n+2}{2}\).
\end{solution}

\section{附加题}

\begin{problem}
求函数 \(y=\e^{-2x}\sin\left(5x+\frac{\pi}{4}\right)\) 的导数.
\end{problem}

\begin{solution}
由乘积法则和复合函数求导法则，得
\[
y'=\left(\e^{-2x}\right)'\sin\left(5x+\frac{\pi}{4}\right)+\e^{-2x}\left[\sin\left(5x+\frac{\pi}{4}\right)\right]'
\]
\[
=-2\e^{-2x}\sin\left(5x+\frac{\pi}{4}\right)+5\e^{-2x}\cos\left(5x+\frac{\pi}{4}\right)
\].
所以
\[
y'=\e^{-2x}\left[5\cos\left(5x+\frac{\pi}{4}\right)-2\sin\left(5x+\frac{\pi}{4}\right)\right]
\].
\end{solution}

\begin{problem}
求定积分：\(\int_0^1\left(x\e^{x^2}+x^2\e^2\right)\mathrm{d}x\).
\end{problem}

\begin{solution}
分别计算两项. 令 \(u=x^2\)，则 \(\mathrm{d}u=2x\mathrm{d}x\)，所以
\[
\int_0^1x\e^{x^2}\mathrm{d}x=\frac{1}{2}\int_0^1\e^u\mathrm{d}u=\frac{\e-1}{2}
\].
又因为 \(\e^2\) 是常数，
\[
\int_0^1x^2\e^2\mathrm{d}x=\e^2\left.\frac{x^3}{3}\right|_0^1=\frac{\e^2}{3}
\].
因此
\[
\int_0^1\left(x\e^{x^2}+x^2\e^2\right)\mathrm{d}x=\frac{\e-1}{2}+\frac{\e^2}{3}=\frac{2\e^2+3\e-3}{6}
\].
\end{solution}
